\documentclass[12pt,a4paper]{article}
\usepackage[utf8]{inputenc}
\usepackage[T1]{fontenc}
\usepackage[brazilian]{babel}
\usepackage{amsmath,amssymb}
\usepackage[margin=2.2cm]{geometry}
\usepackage{tikz}
\usetikzlibrary{arrows.meta,positioning,decorations.pathreplacing}
\usepackage{tcolorbox}
\tcbuselibrary{skins,breakable}
\usepackage{enumitem}
\usepackage{xcolor}
\usepackage{booktabs}
\usepackage{array}
\usepackage{hyperref}
\hypersetup{colorlinks=true,linkcolor=blue!50!black,urlcolor=blue!50!black}

\newtcolorbox{definicao}[1]{colback=blue!5,colframe=blue!55!black,fonttitle=\bfseries,coltitle=white,title={#1},breakable,boxrule=0.8pt,arc=2mm}
\newtcolorbox{exemplo}[1]{colback=green!5,colframe=green!35!black,fonttitle=\bfseries,coltitle=white,title={#1},breakable,boxrule=0.8pt,arc=2mm}
\newtcolorbox{atencao}[1]{colback=red!5,colframe=red!55!black,fonttitle=\bfseries,coltitle=white,title={#1},breakable,boxrule=0.8pt,arc=2mm}
\newtcolorbox{dica}[1]{colback=yellow!12,colframe=orange!70!black,fonttitle=\bfseries,coltitle=black,title={#1},breakable,boxrule=0.8pt,arc=2mm}

\title{\vspace{-1.5cm}\textbf{Tema 4 — Polinômios, Produtos Notáveis e Fatoração}\\[2pt]\large Provavelmente o tema com mais pontos na prova}
\author{}
\date{}

\begin{document}
\maketitle
\thispagestyle{empty}

\section{O que é um polinômio}

\begin{definicao}{Polinômio}
Uma soma de termos do tipo (coeficiente)$\times x^{\text{expoente natural}}$:
\[
p(x) = a_nx^n + a_{n-1}x^{n-1} + \cdots + a_1x + a_0
\]
O maior expoente com coeficiente não nulo é o \textbf{grau} do polinômio.
Ex.: $p(x) = 5x^3-2x^2+x-1$ tem grau $3$.
\end{definicao}

\section{Operações com polinômios}

\subsection{Soma e subtração — só junta termos ``iguais'' (mesmo expoente)}

\begin{exemplo}{$p(x)+q(x)$ com $p(x)=5x^3-2x^2+x-1$ e $q(x)=x^3-x+2$}
Alinhe por grau e some coeficiente com coeficiente:
\[
(5x^3+x^3) + (-2x^2+0) + (x-x) + (-1+2) = 6x^3-2x^2+0x+1 = 6x^3-2x^2+1
\]
\end{exemplo}

\begin{atencao}{Subtração de polinômio: troque TODOS os sinais}
$q(x)-r(x)$ significa somar $q(x)$ com $\bigl(-r(x)\bigr)$ — inverta o sinal
de \textbf{cada} termo de $r(x)$, não só do primeiro.
\end{atencao}

\subsection{Multiplicação — distributiva, todo mundo com todo mundo}

\begin{exemplo}{$p(x)\cdot s(x)$ com $p(x)=5x^3-2x^2+x-1$, $s(x)=x-2$}
Multiplique cada termo de $p$ por cada termo de $s$, depois some os
semelhantes:
\begin{align*}
p(x)\cdot s(x) &= 5x^3\cdot x + 5x^3\cdot(-2) -2x^2\cdot x -2x^2\cdot(-2) + x\cdot x + x\cdot(-2) -1\cdot x -1\cdot(-2)\\
&= 5x^4 -10x^3 -2x^3+4x^2 +x^2-2x -x+2\\
&= 5x^4 -12x^3+5x^2-3x+2
\end{align*}
\end{exemplo}

\subsection{Divisão — método da chave}

\begin{dica}{Como funciona a divisão longa}
Ideia: em cada passo, elimine o termo de \textbf{maior} grau que sobrou,
multiplicando o divisor por um monômio adequado e subtraindo. Repita até o
resto ter grau menor que o divisor.
\end{dica}

\begin{exemplo}{$(x^3+2x^2-4x+1)\div(x-1)$}
\begin{center}
\renewcommand{\arraystretch}{1.3}
\begin{tabular}{r}
$x^3+2x^2-4x+1 \ \big|\ x-1$\\
\hline
\end{tabular}
\end{center}
\textbf{Passo 1}: $x^3 \div x = x^2$. Multiplique $x^2(x-1)=x^3-x^2$ e subtraia:
\[
(x^3+2x^2-4x+1) - (x^3-x^2) = 3x^2-4x+1
\]
\textbf{Passo 2}: $3x^2\div x = 3x$. Multiplique $3x(x-1)=3x^2-3x$ e subtraia:
\[
(3x^2-4x+1)-(3x^2-3x) = -x+1
\]
\textbf{Passo 3}: $-x\div x=-1$. Multiplique $-1(x-1)=-x+1$ e subtraia:
\[
(-x+1)-(-x+1) = 0
\]
Quociente: $x^2+3x-1$. Resto: $0$ (divisão exata).
\[
x^3+2x^2-4x+1 = (x-1)(x^2+3x-1)
\]
\end{exemplo}

\subsection{Dispositivo de Briot--Ruffini (atalho quando o divisor é $x-a$)}

\begin{definicao}{Como montar o dispositivo}
Só funciona quando o divisor tem a forma $x-a$ (grau 1, coeficiente de $x$
igual a 1). Escreva os coeficientes de $p(x)$ em ordem decrescente de grau
(zero para termo ausente) e o valor de $a$ à esquerda:
\begin{center}
\begin{tabular}{c|cccc}
 & $c_3$ & $c_2$ & $c_1$ & $c_0$\\
\hline
$a$ & $c_3$ & $a c_3+c_2$ & $a(\cdot)+c_1$ & $a(\cdot)+c_0$
\end{tabular}
\end{center}
Regra: \textbf{abaixe} o primeiro coeficiente; a partir daí, cada novo
número é (número anterior $\times\, a$) $+$ coeficiente da coluna. Os
números até o penúltimo formam o quociente (grau uma unidade a menos); o
último número é o \textbf{resto}.
\end{definicao}

\begin{exemplo}{$(x^3+2x^2-4x+1)\div(x-1)$ por Briot--Ruffini}
Aqui $a=1$, coeficientes $1,\,2,\,-4,\,1$:
\begin{center}
\begin{tabular}{c|cccc}
 & $1$ & $2$ & $-4$ & $1$\\
\hline
$1$ & $1$ & $1\cdot1+2=3$ & $1\cdot3+(-4)=-1$ & $1\cdot(-1)+1=0$
\end{tabular}
\end{center}
Quociente: $x^2+3x-1$, resto $0$ — bate exatamente com a divisão pela
chave acima.
\end{exemplo}

\begin{exemplo}{Reconstruindo $P(x)$ a partir da tabela (tipo de questão da lista)}
Dado o dispositivo
\begin{center}
\begin{tabular}{c|ccccc}
 & $3$ & $-4$ & $5$ & $d$ & $e$\\
\hline
$a$ & $b$ & $-10$ & $c$ & $24$ & $40$
\end{tabular}
\end{center}
sabemos que $b=3$ (o primeiro sempre ``abaixa'' igual). Da coluna seguinte:
$a\cdot 3+(-4)=-10 \Rightarrow a=-2$. Continuando a regra em cada coluna, dá
para achar $c$, $d$, $e$ um por um, e por fim montar
$P(x)=3x^4-4x^3+5x^2+dx+e$ com os valores de $d,e$ que você calculou.
\textbf{A lição}: cada célula da tabela obedece sempre a mesma conta —
memorize o \emph{processo}, não valores.
\end{exemplo}

\section{Produtos notáveis — decore o \emph{porquê}, não a fórmula}

\begin{exemplo}{Quadrado da soma: $(a+b)^2 = a^2+2ab+b^2$ — visualizando como área}
\begin{center}
\begin{tikzpicture}[scale=0.9]
  \draw[thick] (0,0) rectangle (5,5);
  \draw[thick] (3,0) -- (3,5);
  \draw[thick] (0,3) -- (5,3);
  \node at (1.5,4) {$a^2$};
  \node at (4,4) {$ab$};
  \node at (1.5,1.5) {$ab$};
  \node at (4,1.5) {$b^2$};
  \node[below] at (1.5,0) {$a$};
  \node[below] at (4,0) {$b$};
  \node[left] at (0,4) {$a$};
  \node[left] at (0,1.5) {$b$};
  \draw[decorate,decoration={brace,amplitude=5pt},thick] (0,5.3) -- (5,5.3);
  \node[above] at (2.5,5.6) {$a+b$};
\end{tikzpicture}
\end{center}
A área total do quadrado grande é $(a+b)^2$. Mas somando as 4 partes:
$a^2 + ab+ab+b^2 = a^2+2ab+b^2$. As duas contas descrevem a \textbf{mesma}
área — por isso a fórmula é verdadeira, não é coincidência.
\end{exemplo}

\begin{center}
\renewcommand{\arraystretch}{1.7}
\begin{tabular}{@{}ll@{}}
\toprule
\textbf{Nome} & \textbf{Fórmula}\\
\midrule
Quadrado da soma & $(a+b)^2 = a^2+2ab+b^2$\\
Quadrado da diferença & $(a-b)^2 = a^2-2ab+b^2$\\
Produto da soma pela diferença & $(a+b)(a-b) = a^2-b^2$\\
Cubo da soma & $(a+b)^3 = a^3+3a^2b+3ab^2+b^3$\\
Cubo da diferença & $(a-b)^3 = a^3-3a^2b+3ab^2-b^3$\\
\bottomrule
\end{tabular}
\end{center}

\begin{atencao}{O erro mais famoso da matemática do ensino médio}
$(a+b)^2 \neq a^2+b^2$! Falta o termo do meio $2ab$. Veja pelo desenho
acima: o quadrado grande tem 4 pedaços, não só 2.
\end{atencao}

\begin{exemplo}{Aplicando (estilo da lista)}
\[
\left(\frac{x}{2}+\frac{y}{4}\right)^2 = \frac{x^2}{4} + 2\cdot\frac{x}{2}\cdot\frac{y}{4} + \frac{y^2}{16} = \frac{x^2}{4}+\frac{xy}{4}+\frac{y^2}{16}
\]
\[
\left(\frac{c}{3}+\frac{d}{2}\right)\left(\frac{c}{3}-\frac{d}{2}\right) = \frac{c^2}{9}-\frac{d^2}{4}
\]
\[
(1+a)^3 = 1+3a+3a^2+a^3
\]
\end{exemplo}

\section{Fatoração — o caminho de volta}

\begin{dica}{A pergunta certa a se fazer}
Fatorar é escrever uma soma como um \textbf{produto}. Sempre pergunte, nesta
ordem:
\begin{enumerate}[nosep]
  \item Existe fator comum em todos os termos? (coloque em evidência)
  \item Tem 4 termos? Tente agrupar 2 a 2.
  \item É diferença de quadrados ($a^2-b^2$)?
  \item É trinômio quadrado perfeito ($a^2\pm2ab+b^2$)?
  \item É soma/diferença de cubos?
  \item É um trinômio do 2º grau qualquer? Use Bhaskara para achar as raízes.
\end{enumerate}
\end{dica}

\begin{exemplo}{1. Fator comum em evidência}
\[
15x^4-10x^3 = 5x^3(3x-2)
\]
\end{exemplo}

\begin{exemplo}{2. Agrupamento}
\[
x^3+4x^2+3x+12 = x^2(x+4)+3(x+4) = (x+4)(x^2+3)
\]
\end{exemplo}

\begin{exemplo}{3. Diferença de quadrados}
\[
x^2-9 = (x-3)(x+3) \qquad\qquad x^4-16 = (x^2-4)(x^2+4) = (x-2)(x+2)(x^2+4)
\]
\end{exemplo}

\begin{exemplo}{4. Trinômio quadrado perfeito}
Reconheça: primeiro e último termos são quadrados, e o termo do meio é
$2\times$ raiz do primeiro $\times$ raiz do último.
\[
x^2+10x+25 = (x+5)^2 \qquad (\text{pois } 2\cdot x\cdot5=10x)
\]
\[
16x^2-24xy+9y^2 = (4x-3y)^2 \qquad (\text{pois } 2\cdot4x\cdot3y=24xy)
\]
\end{exemplo}

\begin{exemplo}{5. Soma e diferença de cubos}
\[
a^3-b^3 = (a-b)(a^2+ab+b^2) \qquad\quad a^3+b^3=(a+b)(a^2-ab+b^2)
\]
\[
x^3-27 = (x-3)(x^2+3x+9) \qquad\qquad x^3+64 = (x+4)(x^2-4x+16)
\]
\end{exemplo}

\begin{exemplo}{6. Trinômio do 2º grau geral $ax^2+bx+c$, via Bhaskara}
Ache as raízes $x_1,x_2$ (fórmula na próxima apostila, Tema Equações) e
fatore como $a(x-x_1)(x-x_2)$:
\[
x^2-9x+14: \quad x_1=2,\ x_2=7 \ \Rightarrow\ x^2-9x+14=(x-2)(x-7)
\]
\[
x^2+x-6: \quad x_1=2,\ x_2=-3 \ \Rightarrow\ x^2+x-6 = (x-2)(x+3)
\]
\end{exemplo}

\section{Simplificando frações algébricas — o motivo de fatorar}

\begin{dica}{Regra de ouro}
Só se pode ``cancelar'' fatores que multiplicam o numerador e o
denominador \textbf{inteiros} — nunca corte um termo de dentro de uma soma.
Por isso fatoração vem sempre antes de simplificar.
\end{dica}

\begin{exemplo}{$\dfrac{x^2-9}{x-3}$}
\[
\frac{x^2-9}{x-3} = \frac{(x-3)(x+3)}{x-3} = x+3 \qquad (x\neq3)
\]
\end{exemplo}

\begin{exemplo}{$\dfrac{x^2-4x+4}{x^2-4}$}
\[
\frac{x^2-4x+4}{x^2-4} = \frac{(x-2)^2}{(x-2)(x+2)} = \frac{x-2}{x+2} \qquad (x\neq\pm2)
\]
\end{exemplo}

\begin{atencao}{Cancelamento inválido — clássico da prova}
\[
\frac{x^2+3}{x} \ \neq\ x+3
\]
Você \textbf{não} pode cortar o $x^2$ com o $x$ de baixo, porque $x^2$ está
somado com $3$ (não multiplicado). Só cancela fator de produto completo.
\end{atencao}

\section*{Resumo relâmpago}
\begin{itemize}[nosep]
  \item Soma/subtração: junte só termos de mesmo expoente.
  \item Divisão de polinômio por $(x-a)$: use Briot--Ruffini (mais rápido que a chave).
  \item Produtos notáveis existem para você \textbf{reconhecer padrões}, não para expandir na mão toda vez.
  \item Fatorar = escrever como produto; siga a ordem: comum $\to$ agrupamento $\to$ diferença de quadrados $\to$ trinômio quadrado perfeito $\to$ soma/diferença de cubos $\to$ Bhaskara.
  \item Só simplifique frações depois de fatorar tudo.
\end{itemize}

\end{document}
